% Options for packages loaded elsewhere
\PassOptionsToPackage{unicode}{hyperref}
\PassOptionsToPackage{hyphens}{url}
%
\documentclass[
]{article}
\usepackage{amsmath,amssymb}
\usepackage{iftex}
\ifPDFTeX
  \usepackage[T1]{fontenc}
  \usepackage[utf8]{inputenc}
  \usepackage{textcomp} % provide euro and other symbols
\else % if luatex or xetex
  \usepackage{unicode-math} % this also loads fontspec
  \defaultfontfeatures{Scale=MatchLowercase}
  \defaultfontfeatures[\rmfamily]{Ligatures=TeX,Scale=1}
\fi
\usepackage{lmodern}
\ifPDFTeX\else
  % xetex/luatex font selection
\fi
% Use upquote if available, for straight quotes in verbatim environments
\IfFileExists{upquote.sty}{\usepackage{upquote}}{}
\IfFileExists{microtype.sty}{% use microtype if available
  \usepackage[]{microtype}
  \UseMicrotypeSet[protrusion]{basicmath} % disable protrusion for tt fonts
}{}
\makeatletter
\@ifundefined{KOMAClassName}{% if non-KOMA class
  \IfFileExists{parskip.sty}{%
    \usepackage{parskip}
  }{% else
    \setlength{\parindent}{0pt}
    \setlength{\parskip}{6pt plus 2pt minus 1pt}}
}{% if KOMA class
  \KOMAoptions{parskip=half}}
\makeatother
\usepackage{xcolor}
\usepackage{color}
\usepackage{fancyvrb}
\newcommand{\VerbBar}{|}
\newcommand{\VERB}{\Verb[commandchars=\\\{\}]}
\DefineVerbatimEnvironment{Highlighting}{Verbatim}{commandchars=\\\{\}}
% Add ',fontsize=\small' for more characters per line
\newenvironment{Shaded}{}{}
\newcommand{\AlertTok}[1]{\textcolor[rgb]{1.00,0.00,0.00}{\textbf{#1}}}
\newcommand{\AnnotationTok}[1]{\textcolor[rgb]{0.38,0.63,0.69}{\textbf{\textit{#1}}}}
\newcommand{\AttributeTok}[1]{\textcolor[rgb]{0.49,0.56,0.16}{#1}}
\newcommand{\BaseNTok}[1]{\textcolor[rgb]{0.25,0.63,0.44}{#1}}
\newcommand{\BuiltInTok}[1]{\textcolor[rgb]{0.00,0.50,0.00}{#1}}
\newcommand{\CharTok}[1]{\textcolor[rgb]{0.25,0.44,0.63}{#1}}
\newcommand{\CommentTok}[1]{\textcolor[rgb]{0.38,0.63,0.69}{\textit{#1}}}
\newcommand{\CommentVarTok}[1]{\textcolor[rgb]{0.38,0.63,0.69}{\textbf{\textit{#1}}}}
\newcommand{\ConstantTok}[1]{\textcolor[rgb]{0.53,0.00,0.00}{#1}}
\newcommand{\ControlFlowTok}[1]{\textcolor[rgb]{0.00,0.44,0.13}{\textbf{#1}}}
\newcommand{\DataTypeTok}[1]{\textcolor[rgb]{0.56,0.13,0.00}{#1}}
\newcommand{\DecValTok}[1]{\textcolor[rgb]{0.25,0.63,0.44}{#1}}
\newcommand{\DocumentationTok}[1]{\textcolor[rgb]{0.73,0.13,0.13}{\textit{#1}}}
\newcommand{\ErrorTok}[1]{\textcolor[rgb]{1.00,0.00,0.00}{\textbf{#1}}}
\newcommand{\ExtensionTok}[1]{#1}
\newcommand{\FloatTok}[1]{\textcolor[rgb]{0.25,0.63,0.44}{#1}}
\newcommand{\FunctionTok}[1]{\textcolor[rgb]{0.02,0.16,0.49}{#1}}
\newcommand{\ImportTok}[1]{\textcolor[rgb]{0.00,0.50,0.00}{\textbf{#1}}}
\newcommand{\InformationTok}[1]{\textcolor[rgb]{0.38,0.63,0.69}{\textbf{\textit{#1}}}}
\newcommand{\KeywordTok}[1]{\textcolor[rgb]{0.00,0.44,0.13}{\textbf{#1}}}
\newcommand{\NormalTok}[1]{#1}
\newcommand{\OperatorTok}[1]{\textcolor[rgb]{0.40,0.40,0.40}{#1}}
\newcommand{\OtherTok}[1]{\textcolor[rgb]{0.00,0.44,0.13}{#1}}
\newcommand{\PreprocessorTok}[1]{\textcolor[rgb]{0.74,0.48,0.00}{#1}}
\newcommand{\RegionMarkerTok}[1]{#1}
\newcommand{\SpecialCharTok}[1]{\textcolor[rgb]{0.25,0.44,0.63}{#1}}
\newcommand{\SpecialStringTok}[1]{\textcolor[rgb]{0.73,0.40,0.53}{#1}}
\newcommand{\StringTok}[1]{\textcolor[rgb]{0.25,0.44,0.63}{#1}}
\newcommand{\VariableTok}[1]{\textcolor[rgb]{0.10,0.09,0.49}{#1}}
\newcommand{\VerbatimStringTok}[1]{\textcolor[rgb]{0.25,0.44,0.63}{#1}}
\newcommand{\WarningTok}[1]{\textcolor[rgb]{0.38,0.63,0.69}{\textbf{\textit{#1}}}}
\usepackage{longtable,booktabs,array}
\usepackage{calc} % for calculating minipage widths
% Correct order of tables after \paragraph or \subparagraph
\usepackage{etoolbox}
\makeatletter
\patchcmd\longtable{\par}{\if@noskipsec\mbox{}\fi\par}{}{}
\makeatother
% Allow footnotes in longtable head/foot
\IfFileExists{footnotehyper.sty}{\usepackage{footnotehyper}}{\usepackage{footnote}}
\makesavenoteenv{longtable}
\setlength{\emergencystretch}{3em} % prevent overfull lines
\providecommand{\tightlist}{%
  \setlength{\itemsep}{0pt}\setlength{\parskip}{0pt}}
\setcounter{secnumdepth}{-\maxdimen} % remove section numbering

\usepackage{newunicodechar}
\usepackage{amssymb}
\newunicodechar{λ}{\ensuremath{\lambda}}
\newunicodechar{κ}{\ensuremath{\kappa}}
\newunicodechar{σ}{\ensuremath{\sigma}}
\newunicodechar{₁}{\ensuremath{_1}}
\newunicodechar{‖}{\ensuremath{\Vert}}
\newunicodechar{≤}{\ensuremath{\le}}
\newunicodechar{→}{\ensuremath{\rightarrow}}
\newunicodechar{·}{\ensuremath{\cdot}}
\newunicodechar{ℚ}{\ensuremath{\mathbb{Q}}}
\newunicodechar{∈}{\ensuremath{\in}}
\newunicodechar{…}{\ldots}
\newunicodechar{§}{\S}
\newunicodechar{—}{---}
\newunicodechar{–}{--}
\newunicodechar{’}{'}
\newunicodechar{“}{``}
\newunicodechar{”}{''}
\newunicodechar{⁠}{}
\ifLuaTeX
  \usepackage{selnolig}  % disable illegal ligatures
\fi
\IfFileExists{bookmark.sty}{\usepackage{bookmark}}{\usepackage{hyperref}}
\IfFileExists{xurl.sty}{\usepackage{xurl}}{} % add URL line breaks if available
\urlstyle{same}
\hypersetup{
  hidelinks,
  pdfcreator={LaTeX via pandoc}}

\author{}
\date{}

\begin{document}

\hypertarget{pm-ll-pirtm-004-the-three-hashes-witness-address-and-external-seal}{%
\section{PM-LL-PIRTM-004 --- The Three Hashes: Witness, Address, and
External
Seal}\label{pm-ll-pirtm-004-the-three-hashes-witness-address-and-external-seal}}

Status: architectural demarcation. Date: 2026-10-12 Published anchor:
PhaseMirror/phase-mirror-agent-dev
\texttt{48a1752d08206a0baff1d22e33834309fd85e5b5}

\hypertarget{boundary-statement}{%
\subsection{Boundary Statement}\label{boundary-statement}}

The stratification witness
(\texttt{crates/lexlean-stratum/tests/stratification\_witness.rs}) binds
a LexLean receipt to its own domain. It does not bind, recompute, or
verify a PIRTM seal.

\hypertarget{what-the-witness-binds}{%
\subsection{What the Witness Binds}\label{what-the-witness-binds}}

At \texttt{48a1752} the witness executes:

\begin{Shaded}
\begin{Highlighting}[]
\PreprocessorTok{assert\_eq!}\NormalTok{(}\PreprocessorTok{ContractivityReceipt::}\NormalTok{HASH\_AUTHORITY}\OperatorTok{,} \StringTok{"lexlean{-}contractivity{-}v1"}\NormalTok{)}\OperatorTok{;}
\PreprocessorTok{assert\_eq!}\NormalTok{(}
\NormalTok{    receipt\_fields}\OperatorTok{.}\NormalTok{get(}\StringTok{"hash\_authority"}\NormalTok{)}\OperatorTok{,}
    \ConstantTok{Some}\NormalTok{(}\OperatorTok{\&}\PreprocessorTok{Json::}\BuiltInTok{Str}\NormalTok{(}\StringTok{"lexlean{-}contractivity{-}v1"}\OperatorTok{.}\NormalTok{to\_string()))}
\NormalTok{)}\OperatorTok{;}
\end{Highlighting}
\end{Shaded}

\texttt{HASH\_AUTHORITY} is a constant on the receipt. The second check
guarantees that the canonical JSON carried by the entry names that
authority. That is the whole binding: it names the authority of the
hash. It does not hash anything against a PIRTM value.

The receipt hash is SHA-256 of a §21.1 frame encoding under domain
\texttt{lexlean-contractivity-v1}, over six frames: -
\texttt{prime-index} - \texttt{atom-count} - \texttt{anchor} -
\texttt{norm-numerator} - \texttt{norm-denominator} - \texttt{status}

For the fixture
\texttt{def\ a\ :\ Nat\ :=\ 1\textbackslash{}ndef\ b\ :\ Nat\ :=\ a\textbackslash{}n},
that is prime \texttt{2}, two atoms, anchor \texttt{a}, norm
\texttt{1/6}, status \texttt{ACCEPT}. SPEC.md §34.2 and \texttt{LP-08}
require the entry to state that this is not the compiler's
\texttt{seal\_hash}.

SnapAddr is a second LexLean identity under a distinct domain,
\texttt{lexlean-pirtm-v1}. Its frames are the canonicalization
identifier, prime index, atom count, and each atom's name and body. For
the fixture, the snapaddr is \texttt{6b16a39b…}. It is an address of the
stratum, not a seal of a contractivity receipt.

\hypertarget{what-the-seal-is-and-what-it-is-not}{%
\subsection{What the Seal Is, and What It Is
Not}\label{what-the-seal-is-and-what-it-is-not}}

SPEC.md §34.0 item 3 is the governing invariant: the compiler's
\texttt{seal\_hash} preimage is not recoverable from the binary that was
probed. A LexLean receipt does not reproduce it and does not claim to.
The entry carries \texttt{hash\_authority} so that nobody looks a
LexLean receipt up in a PIRTM ledger. Computing that seal inside
\texttt{lexlean-stratum} would be the substitution §34.4 forbids.

The public \texttt{SNAPKITTYWEST/sovereign-pirtm} tree is a distinct
artifact and not the preimage of the MCP seal. Its
\texttt{ReceiptGenerator} builds a SHA-256 over a string including
operator name, prime, \texttt{now\_ns()}, and an incrementing counter. A
hash that incorporates wall-clock time and a local counter is not
reproducible. Its \texttt{verify} checks the prime, operator name, and
that the hash length is 64 characters. The README's
\texttt{seal(input,\ output)} and JSON fields
\texttt{PIRTM-CONTRACTIVITY-1} / \texttt{sha256-merkle-v1} do not match
that C++.

\hypertarget{the-three-hashes}{%
\subsection{The Three Hashes}\label{the-three-hashes}}

\begin{longtable}[]{@{}
  >{\raggedright\arraybackslash}p{(\columnwidth - 6\tabcolsep) * \real{0.2500}}
  >{\raggedright\arraybackslash}p{(\columnwidth - 6\tabcolsep) * \real{0.2500}}
  >{\raggedright\arraybackslash}p{(\columnwidth - 6\tabcolsep) * \real{0.2500}}
  >{\raggedright\arraybackslash}p{(\columnwidth - 6\tabcolsep) * \real{0.2500}}@{}}
\toprule\noalign{}
\begin{minipage}[b]{\linewidth}\raggedright
Hash
\end{minipage} & \begin{minipage}[b]{\linewidth}\raggedright
Domain
\end{minipage} & \begin{minipage}[b]{\linewidth}\raggedright
Recomputable from the lexeme
\end{minipage} & \begin{minipage}[b]{\linewidth}\raggedright
What it certifies
\end{minipage} \\
\midrule\noalign{}
\endhead
\bottomrule\noalign{}
\endlastfoot
LexLean receipt & \texttt{lexlean-contractivity-v1} & Yes & \(\|G\|_1\)
exact rational norm of the declaration graph \\
SnapAddr & \texttt{lexlean-pirtm-v1} & Yes & Identity of the stratum \\
PIRTM \texttt{seal\_hash} & Not published for the probed binary & No &
Whatever that binary sealed; not claimed here \\
\end{longtable}

\hypertarget{process-boundary-quarantine}{%
\subsection{Process-Boundary
Quarantine}\label{process-boundary-quarantine}}

A process-boundary call to \texttt{pirtm\_verify\_ensemble} can store
the returned seal as external evidence, labeled as PIRTM's (Option 3
pattern). It cannot become the LexLean receipt, and the crate cannot
grow a function that pretends to produce it.

\end{document}
